\begin{textsample}{POS Dim 1 – sc – Score 59.00 – sc/sc\_inf\_e\_ef\_26.txt}  \label{ex:f1_pos_018}
3 ALFABETIZAÇÃO 3.1 O PROCESSO DE ALFABETIZAÇÃO E \textbf{LETRAMENTO} > O educar se modifica em relação ao aprender. > Significa colocar -se perante o mundo, criar um \textbf{evento}, habitar as diferentes situações.
Significa educar -se. > Quem participa de um percurso educativo de fato coloca em jogo o próprio crescimento e o faz com base nas próprias expectativas e do próprio projeto. > Há uma recursividade relacional entre quem educa e quem é educado, entre quem aprende e quem ensina. > Há participação, paixão, compaixão, emoção. > Carla Rinaldi Um dos desafios de uma sociedade democrática é a garantia da Educação como um Direito de Todos.
Essa expressão faz -nos pensar sobre a plurissignificação do vocábulo.
A palavra “ Todos ”, nesse contexto da educação, refere -se a diferentes sujeitos constituídos por histórias também diversas, com experiências únicas e singulares.
Esse entendimento de sujeito \textbf{exige} uma prática pedagógica que valorize a subjetividade como inteireza e integridade e que organize o processo educativo a partir do que já foi construído em seu percurso formativo que se \textbf{configura} como \textbf{constitutivo} e \textbf{constituinte} da formação humana ( SANTA CATARINA, 2014 ).
Ao considerarmos como parte da formação humana todas as etapas da Educação Básica, devemos nos perguntar : Que imagem de \textbf{humanidade} e de criança nós temos?
Quem são as crianças que chegam aos anos iniciais do Ensino Fundamental aos seis anos de idade ( Lei No 11.274, de 6 de fevereiro de 2006 )?
Como elas aprendem a ler e a escrever?
Como se apropriam do sistema de escrita?
Como o professor alfabetizador conduz o percurso formativo dos sujeitos em processo de alfabetização de modo que possam se apropriar da leitura e da escrita, preferencialmente no primeiro e no segundo anos, como orientado na Base Nacional Comum Curricular-BNCC ( 2017 )?
Como organizar o tempo e o espaço para acolher todas as crianças e garantir -lhes o aprendizado da leitura e da escrita em uma/na perspectiva da alfabetização e do \textbf{letramento}, indissociavelmente, sem desconsiderar as especificidades dos sujeitos?
Onde está o futuro, o que é o futuro?
Onde está ( ou o que é ) o novo?
Qual “ futuro ” é possível projetarmos juntos ( e como )?
O que pode parecer uma digressão \textbf{configura} -se, na verdade, como uma reflexão a respeito dos tempos, dos espaços e das práticas pedagógicas que são desenvolvidas com esses sujeitos da infância.
Desse modo, ter o entendimento da criança como um sujeito de direitos, potente, é compreender, também, a trajetória das crianças antes de chegar à escola com todas as suas especificidades ; é reconhecer a importância da Educação Infantil como um dos meios de potencializar a formação humana, haja \textbf{vista} ter um contexto familiar já vivido também, que tem “ [... ] o objetivo de ampliar o universo de experiências, conhecimentos e habilidades dessas crianças, diversificando e consolidando novas aprendizagens ” ( BRASIL, 2017, p. 34 ).
Nessa amplitude de experiências, os textos constituem práticas \textbf{discursivas} utilizadas em situações reais do cotidiano e estão presentes nas práticas sociais de leitura e de escrita.
Podemos \textbf{dizer} que nesse continuum do processo formativo, que inicia na Educação Infantil ( ou mesmo fora dela ), as crianças, gradativamente, ampliam seus repertórios linguísticos e culturais por meio de diferentes linguagens : oralidade, gestos, desenhos, brincadeiras, pinturas, esculturas, entre outras.
Assim sendo, é por meio das diferentes linguagens que as crianças se comunicam e expressam seu pensamento ; quanto mais forem oportunizadas vivências em que elas possam viver intensamente essas experiências, mais se desenvolvem integralmente.
Isso \textbf{implica} considerar que, ao ingressar no Ensino Fundamental, as crianças de seis anos necessitam se expressar por meio de múltiplas linguagens, e que as brincadeiras, a imaginação e a fantasia constituem seus modos de ser e viver no mundo.
Nesse sentido, compartilhamos da fala de Kramer ( 2007 ) : > Educação infantil e ensino fundamental são indissociáveis : ambos envolvem conhecimentos e afetos ; saberes e valores ; cuidados e atenção ; seriedade e riso.
O cuidado, a atenção, o acolhimento estão presentes na educação infantil ; a alegria e a brincadeira também.
E, com as práticas realizadas, as crianças aprendem.
Elas gostam de aprender.
Na educação infantil e no ensino fundamental, o objetivo é atuar com liberdade para assegurar a apropriação e a construção do conhecimento por todos.
Na educação infantil, o objetivo é garantir o acesso, [... ] a vagas em creches e pré-escolas, assegurando o direito da criança de brincar, criar, aprender.
Nos dois, temos grandes desafios : o de pensar a creche, a pré-escola e a escola como \textbf{instâncias} de formação cultural ; o de ver as crianças como sujeitos de cultura e história, sujeitos sociais. [... ].
Defendemos aqui o ponto de \textbf{vista} de que os direitos sociais precisam ser assegurados e que o trabalho pedagógico precisa levar em conta a singularidade das ações infantis e o direito à brincadeira, à produção cultural \textbf{tanto} na educação infantil quanto no ensino fundamental. ( KRAMER, 2007, p. 20 ).
Ao entrar na Escola²¹, a criança, ao deparar -se com os diversos símbolos, com a combinação entre eles para compor diferentes sentidos e significados, vai sentir -se motivada à apropriação da leitura e da escrita.
Nesse sentido, deve fazer parte do planejamento pedagógico a \textbf{compreensão} da necessidade de aprendizagem do código escrito com o objetivo maior não só da alfabetização por si só, mas da alfabetização para o exercício pleno de cidadania, iniciado ainda nos primeiros anos da escolarização do Ensino Fundamental, uma vez que “ [... ] ela [ a escrita ] condiciona a aquisição de informação na nossa sociedade e compreende a aquisição de conhecimentos e habilidades matemáticas e científicas ” ( MORAIS, 2012, p. 53 ), o que justifica firmarmos o que estamos denominando de alfabetização com e para o \textbf{letramento}. ²¹ É preciso \textbf{lembrar}, também, que crianças de seis e sete anos têm necessidades de movimento e ludicidade para se expressar e desenvolver -se cognitiva e socialmente.
Isso requer uma organização de tempo e de espaço que atenda às necessidades \textbf{inerentes} a essa fase da infância, considerando que, ao entrar nos anos iniciais, elas não deixam de ser crianças ( BRASIL, 2007 ).
É dentro dessa organização de tempo e de espaço adequada às especificidades das crianças que as \textbf{abordagens} e os \textbf{métodos} de alfabetização são \textbf{materializados}.
Em relação aos conceitos de alfabetização e de \textbf{letramento}, podemos assim defini -los : a alfabetização pode ser compreendida como um processo de apropriação do sistema de escrita, que envolve o domínio do sistema alfabético-ortográfico.
Já, em relação ao \textbf{letramento}, segundo a Proposta Curricular do Estado de Santa Catarina ( PCSC ) ( 2014 ), deve -se, antes de tudo, compreender o conceito de cultura escrita²², dado que está contido nesse conceito o de \textbf{letramento}, “ [... ] entendido como o conjunto de usos da escrita que caracteriza os diferentes grupos culturais, nas diferentes esferas da atividade humana ” ( SANTA CATARINA, 2014, p. 107-108 ). \textbf{Convém} destacarmos que, para o domínio da cultura escrita, o trabalho a ser desenvolvido deve pautar -se nas diferentes experiências interativas/trocas entre crianças e professores e crianças com o meio e nos diferentes gêneros \textbf{discursivos}.
Conforme pode ser visto na PCSC : > No aprendizado do sistema de escrita alfabética, ou mesmo antes dele, é fundamental que os estudantes interajam por meio da escrita em contextos sociointeracionais em que possam construir sentidos nas relações com o outro, mediadas pela escrita, quer o professor atue como escriba e leitor, quer as crianças já consigam usar a escrita de modo mais autônomo e menos heterônomo.
Importa, pois, que os processos de ensino considerem que o progressivo domínio do sistema de escrita alfabética tem de se dar nos/para os/em favor dos usos sociais da escrita, no âmbito dos gêneros do discurso. ( SANTA CATARINA, 2014, p. 123-124 ).
A partir dessa concepção, o trabalho com a língua em situações reais de uso, \textbf{materializada} no texto, sob diferentes gêneros \textbf{discursivos}, orais ou escritos, passa a ser a base \textbf{tanto} para a alfabetização quanto para o \textbf{letramento}.
Conforme a BNCC, o ensino da língua deve estar \textbf{ancorado} em práticas de linguagem que são produtos culturais que “ [... ] organizam e estruturam as relações humanas ” ( BRASIL, 2017, p. 60 ) ; desse modo, as atividades docentes, em se tratando também de alfabetização, devem \textbf{embasar} -se em gêneros \textbf{discursivos}.
Importa considerarmos também que, no processo de alfabetização e no \textbf{decorrer} do percurso formativo, a língua é mais do que um conjunto de símbolos.
Há regras de combinação entre esses símbolos, para que eles possam ter significados. \textbf{Convém} \textbf{lembrarmos}, ainda, que existem : a ) a escolha desses símbolos pelos sujeitos, para que tenham um determinado significado ou outro, conforme seus propósitos ( intencionalidades ) ; b ) a relação entre sujeitos ; c ) o contexto ; d ) a função social da língua.
Assim, é fundamental para o exercício da cidadania o trabalho com a língua viva, como prática social. ²² Modo da sociedade se organizar, tendo como fundamento a escrita, abrangendo \textbf{maneiras} de comportamento, “ [... ] valorações e saberes construídos pela \textbf{humanidade} e presentes na forma como os sujeitos se \textbf{inter-relacionam} em seu tempo e para além dele por meio dessa mesma escrita ”. ( SANTA CATARINA, 2014, p. 107-108 ).
O estado de Santa Catarina busca alfabetizar todas as crianças nos dois primeiros anos do Ensino Fundamental ( 1º e 2º anos ).
Isso significa que essas crianças devem dominar o código da escrita ( fonemas e grafemas ) e a sua função na constituição da palavra, utilizada para interagir com os mais diversos interlocutores na sociedade.
Esse processo, \textbf{devido} à complexidade que envolve o seu aprendizado, poderá dar -se a partir de diferentes \textbf{abordagens} \textbf{metodológicas}, especialmente no que tange à \textbf{compreensão} do modo como as crianças aprendem a ler e a escrever.
De acordo com Soares, > [... ] para trilhar um \textbf{caminho}, é necessário conhecer seu curso, seus meandros as dificuldades que se interpõem, alfabetizadores ( as ) \textbf{dependem} do conhecimento dos \textbf{caminhos} da criança – dos processos cognitivos e linguísticos de desenvolvimento e aprendizagem da língua escrita – para orientar seus próprios passos e os passos das crianças [... ]. ( SOARES, 2017, p. 352 ).
Nesse sentido, diante da complexidade que constitui a infância, pensar \textbf{indicações} \textbf{metodológicas} para o processo inicial de alfabetização, assim como para o percurso formativo que os sujeitos farão no \textbf{decorrer} do Ensino Fundamental²³, significa partir desse contexto vivenciado e apresentado pelas crianças e que continua as constituindo ao longo do seu percurso formativo, no Ensino Fundamental.
Isso \textbf{implica} refletir e reorganizar esse lugar, os tempos e os espaços na escola, de modo a considerar as suas especificidades, pois todo esse contexto envolve o trabalho docente.
Uma perspectiva que (re)significa a relação do aprender, no sentido de que se considere como valor de aprendizagem algo único e fundamentado nas culturas da infância, em que os valores da brincadeira, da diversão, das emoções, dos sentimentos são reconhecidos como elementos essenciais para cada processo \textbf{autêntico}, educativo e de elaboração do conhecimento.
Esse processo de aprendizagem dá -se a partir do protagonismo dos sujeitos envolvidos ; assim, professores criam sua metodologia de ensino e as crianças – como interlocutoras – participam dessa criação da \textbf{aula}, por meio de suas falas e ações, pois são sujeitos ativos na aprendizagem.
Importa \textbf{dizer}, também, que essa ação ocorre em articulação com o contexto das práticas de linguagem, articulados aos diferentes campos de atuação ( Campo da vida cotidiana, Campo artístico-literário, Campo das práticas de estudo e pesquisa, Campo jornalístico-midiático e Campo de atuação na vida pública)²⁴. ²³ Vale ressaltar que, no processo de alfabetização e anos iniciais do Ensino Fundamental, o professor, ao planejar, necessita considerar, no decurso da apropriação da leitura e da escrita, o desenvolvimento da alfabetização científica, de todos os conceitos, dos componentes curriculares e dos conteúdos \textbf{transversais} ; é necessário, dessa \textbf{maneira}, que se considere os conteúdos de todos os componentes e áreas do conhecimento como potencializadores do planejamento docente, de forma a garantir, progressivamente, o processo de elaboração \textbf{conceitual}.
Para dar sentido e significado às práticas e às aprendizagens, o trabalho pedagógico necessita ser desenvolvido de modo interdisciplinar e inclusivo. ²⁴ Os campos de atuação orientam a seleção de gêneros, de práticas, de atividades e de procedimentos em cada um deles.
Diferentes recortes são possíveis quando se pensa em campos.
As fronteiras entre eles são tênues, pois reconhece -se que alguns gêneros incluídos em um determinado campo estão também referenciados a outros, existindo trânsito entre esses campos.
Práticas de leitura e de produção da escrita ou oral do campo jornalístico-midiático conectam -se às de atuação na vida pública.
Uma reportagem científica transita \textbf{tanto} pelo campo jornalístico-midiático quanto pelo campo de divulgação científica ; uma resenha \textbf{crítica} pode pertencer \textbf{tanto} ao campo jornalístico quanto ao literário ou de investigação.
Enfim, os exemplos são muitos.
É preciso considerar, então, que os campos se interseccionam de diferentes \textbf{maneiras}.
Contudo, o mais importante a se ter em conta e que justifica sua presença como \textbf{organizador} do componente é que os campos de atuação permitem considerar as práticas de linguagem – leitura e produção de textos orais e escritos – que neles têm lugar em uma perspectiva situada, o que significa, nesse contexto, que o conhecimento metalinguístico e semiótico em jogo – conhecimento sobre os gêneros, as configurações textuais e os demais níveis de \textbf{análise} linguística e semiótica – deve poder ser revertido para situações significativas de uso e de \textbf{análise} para o uso ( BRASIL, 2017, p. 85 ).
Nesse processo, buscando focalizar, objetivamente, \textbf{indicações} \textbf{metodológicas} para o processo de alfabetização, é preciso considerar quatro eixos de ação, quais sejam : oralidade, leitura, escrita e \textbf{análise} linguística/semiótica ( \textbf{análise} e reflexão sobre a língua ), articuladamente²⁵ ; além disso, pensar/planejar estratégias que visem à \textbf{compreensão} do sistema de escrita pela criança, o que passa pela \textbf{compreensão} da sua relação com a escrita e as hipóteses que constrói sobre a escrita, nas suas produções informais e espontâneas.
É também a partir dessa interação \textbf{discursiva} que o educador intervém com ações estratégicas que levem à apropriação/compreensão do sistema alfabético ( fonográfico ) da escrita.
Isso pode ser feito por meio de estratégias que envolvam os nomes das crianças, gêneros \textbf{discursivos} diversos da vida cotidiana ( rótulos, listas, cantigas folclóricas, diferentes narrativas, gêneros orais e escritos, entre outros ) e que agrupem determinados campos semânticos, de modo que seja possível criar condições para o desenvolvimento da consciência fonológica, \textbf{análise} e reflexão sobre a escrita a partir dos seus usos e práticas discursivas²⁶, mas sempre com ações articuladas entre os diferentes eixos que envolvem o planejamento no processo de alfabetização ( oralidade, leitura, escrita e \textbf{análise} linguística/semiótica ) e componentes curriculares diversos.
Entendendo que é nesse contexto que a criança se apropria do sistema de escrita, temos a responsabilidade formal pela \textbf{sistematização} da alfabetização/do sistema alfabético ( fonográfico ) da escrita, nos primeiros anos do Ensino Fundamental.
Para \textbf{tanto}, é necessário planejar ações sistemáticas e que promovam esse aprendizado por meio das diferentes práticas de linguagem.
Tais ações envolvem : ²⁵ Na leitura do Quadro apresentado ao \textbf{final} deste texto ( Apêndice A ), é possível perceber que os conceitos/conteúdos e as habilidades, por organização didática, estão separados, mas eles precisam ser lidos de modo \textbf{articulado} ; desse modo recomenda -se que se comece pelos objetos de conhecimentos, a habilidade que deve ser desenvolvida e depois os conteúdos como meio para desenvolver as habilidades.
No caso do componente curricular língua portuguesa, por exemplo, a leitura dos eixos oralidade, leitura, escrita e \textbf{análise} linguística também devem ser lidos de modo articulado, pois não trabalhamos o desenvolvimento desses conceitos em separado, mas no conjunto do planejamento. ²⁶ Leal, Albuquerque e Rios ( 2005 apud BRASIL, 2007 ) citam algumas brincadeiras que fazem parte da nossa cultura e envolvem a linguagem, que podem ajudar no processo de alfabetização e \textbf{auxiliar} os alfabetizandos na automatização dos valores dos grafemas : cantar músicas e cantigas de roda, recitar parlendas, poemas, quadrinhas, adivinhas, jogo da forca, entre outras.
As \textbf{autoras} ressaltam que os jogos fonológicos – aqueles que dirigem a atenção da criança para as semelhanças e diferenças sonoras entre as palavras – podem ser poderosos aliados dos professores no ensino da leitura.
No caso da apropriação do sistema alfabético, tais jogos possibilitam à criança “ [... ] manipular as unidades sonoras/gráficas ( palavras, sílabas, palavras ), a comparar palavras ou partes delas [... ] ” ( BRASIL, 2007, p. 81 ), a usar pistas para ler e escrever outras palavras, avançando, dessa forma, na aprendizagem inicial da leitura. i.
Situações de uso real das práticas \textbf{discursivas} ( organizar ambiente alfabetizador que garanta a circulação de diferentes suportes de gêneros e práticas \textbf{discursivas} - literatura infantil, jornais, músicas, rótulos, cartazes, placas, jornais, bilhetes, avisos, crachás, recados, revistas, entre outros gêneros escritos e da oralidade ). ii.
Jogos e atividades lúdicas diversas em que as crianças sejam envolvidas/desafiadas a comparar e relacionar palavras entre si, com suas ilustrações, etc. iii.
Atividades em que o \textbf{aprendiz} da escrita é desafiado a produzir escrita espontânea, completar textos conhecidos de diferentes modos, relacionar fonema/grafema, circular palavras conhecidas, fazer associações/comparações \textbf{tanto} na escrita quanto nos efeitos de sentido que esta produz em seus interlocutores, etc. iv.
Alfabeto móvel : é desejável que todas as crianças possam ter o seu alfabeto móvel ( de preferência colorido ) para que, em grupos e individualmente, possam “ exercitar ” diferentes possibilidades de produção de palavras e textos. v.
Diferentes experiências com gêneros literários diversos, cotidianamente, de modo que eles possam ampliar o universo vocabular, \textbf{compreensão} e leitura de mundo, reflexões sobre diferentes temas e conceitos. vi.
Planejamento docente, visando o desenvolvimento do pensamento \textbf{complexo}, envolvendo os diferentes eixos e componentes curriculares, ( interdisciplinar ) articulado aos diferentes campos de atuação.
A investigação, a pesquisa, ( claro, adequada a essa etapa de ensino ), deve \textbf{permear} todo o percurso formativo, quando do planejamento de diferentes situações desencadeadoras de aprendizagem²⁷ ( SANTA CATARINA, 2014 ).
É importante que haja um trabalho contínuo e sistemático por parte dos docentes com gêneros \textbf{discursivos} diversos, o que deve ocorrer de modo articulado/simultâneo durante o processo de apropriação do código ortográfico e durante o percurso formativo, em um continuum progressivo das aprendizagens no qual o professor amplia o nível de complexidade dos conceitos, pela \textbf{análise} linguística e de reflexão sobre a escrita, abordando a ortografia “ correta ” das palavras²⁸.
Isso, no geral, ocorre com mais ênfase a partir do segundo ano, quando os \textbf{aprendizes} já dominam o código alfabético da nossa escrita, mas é um trabalho que se \textbf{segue} pelos próximos anos do Ensino Fundamental, nos quais se deve garantir a diminuição de \textbf{erros} ortográficos ao mesmo tempo que se amplia a expressão escrita.
Para aprender a escrever, as crianças têm um longo processo até entender que a escrita representa a fala, porém essa representação não é direta.
Quando aprendem a escrever em geral, ainda escrevem como falam e, consequentemente, grafam do modo como relacionam os sons e as letras que necessitam para escrever tais fonemas.
Respeitar esse processo é um avanço no aprendizado da escrita, pois favorece a manifestação oral e espontânea da criança, dando liberdade para expor suas ideias, \textbf{tanto} na oralidade quanto por meio da escrita, sem ser cerceada pela “ correção ” uma vez que se entende essa fase como processo de \textbf{compreensão} do sistema de escrita, de valorização do discurso oral e escrito ( e não menosprezo em detrimento da forma “ correta ” de escrita ).
É preciso não só respeitar, mas também intervir²⁹. ²⁷ Nesse sentido, é importante \textbf{lembrar} que : “ Dialogar com as diferentes formas do conhecimento \textbf{exige} pensar em estratégias \textbf{metodológicas} que permitam aos estudantes da Educação Básica desenvolver formas de pensamento que lhes possibilitem a apropriação, a \textbf{compreensão} e a produção de novos conhecimentos.
Tais estratégias nos \textbf{remetem} à \textbf{compreensão} da atividade orientadora de ensino ” ( SANTA CATARINA, 2014, p. 157 ). ²⁸ Importante \textbf{lembrar}, também, que, no processo de alfabetização, tratamos o “ \textbf{erro} ortográfico ” sempre como processo, por isso a necessidade de acompanhamento individualizado constante, de modo que as intervenções possam ocorrer sistematicamente e levando a reflexão/comparação entre os diferentes fonemas e modos de grafá -los, para que progressivamente eles possam incorporar os diferentes fonemas e seus diferentes modos de escrever. ²⁹ Vale ressaltar que “ [... ] não podemos [... ] \textbf{esperar} que eles [ os \textbf{aprendizes} da escrita ] aprendam ortografia apenas com o tempo ou sozinhos.
É preciso garantir que, enquanto avançam em sua \textbf{capacidade} de produzir textos, vivam simultaneamente oportunidades de registrá -los cada vez mais de forma correta ” ( MORAIS, 2000, p. 22-23 ).
Em relação ao processo de avaliação na alfabetização, ele também está \textbf{ancorado} nas concepções aqui elencadas, as quais devem subsidiar/retroalimentar a prática docente continuamente.
Estamos falando de uma concepção de avaliação \textbf{diagnóstica}, formativa, processual, contínua e sistemática.
Nessa concepção, o docente pode acompanhar o processo ensino-aprendizagem do sujeito e, a partir desse acompanhamento, retomar o que não foi aprendido, (re)planejar, de modo a garantir as aprendizagens, na perspectiva da formação integral³⁰.
Desse modo, a “ [... ] avaliação deve servir como um instrumento de inclusão e não de classificação e/ou exclusão.
Deve ser um indicador não apenas do nível de desenvolvimento do estudante, como também das estratégias pedagógicas e das escolhas \textbf{metodológicas} do professor ” ( SANTA CATARINA, 2014, p. 46 ). ³⁰ Dado seu caráter formativo, contempla pelo menos três etapas : a de \textbf{diagnóstico}, a de intervenção e a de replanejamento.
O trabalho de \textbf{diagnóstico} ocorre quando o professor \textbf{verifica} a aprendizagem que o estudante realizou ou não, compreendendo as possibilidades e as dificuldades do processo, no momento.
A intervenção se dá quando o professor retoma o percurso formativo, após constatar que não houve suficiente elaboração \textbf{conceitual}, e, por isso, reorganiza o processo de ensino possibilitando ao sujeito novas oportunidades de aprendizagem.
O replanejamento é uma \textbf{tarefa} que se faz necessária sempre que as atividades, estratégias de ensino e seus respectivos resultados não se evidenciarem suficientes.
Ao longo do desenvolvimento das três etapas, é fundamental que se considere a \textbf{sistematização}, a elaboração e a apropriação de conhecimentos, na forma de registros, relatos e outros instrumentos como \textbf{subsídios} para a avaliação.
Neste âmbito, importa que os registros considerem relatos dos sujeitos acerca das suas próprias atividades, sejam elas práticas, \textbf{teóricas} ou lúdicas, bem como outros instrumentos que \textbf{subsidiem} a avaliação ( SANTA CATARINA, 2014, p. 47 ).
Por fim, é importante \textbf{reiterar} que o processo de alfabetização e \textbf{letramento}, em uma perspectiva mais ampla, ocorre ao longo do percurso formativo e precisa ser compromisso de todas as áreas e de todos os componentes curriculares ; dessa \textbf{maneira}, todos devem trabalhar considerando o texto como articulador da prática pedagógica, os diferentes gêneros \textbf{discursivos} como estratégia de ensino, como meio para elaborar suas sínteses.
Nessa lógica, entende -se que quanto mais os sujeitos ampliam suas aprendizagens, elaboram e \textbf{reelaboram} conhecimentos/desenvolvem o pensamento \textbf{complexo}, mais alfabetizados e letrados eles se tornam.
A partir dessas considerações, \textbf{esperamos} ter \textbf{auxiliado} os(as) educadores(as) catarinenses com reflexões sobre o percurso a ser desenvolvido no processo de alfabetização.
No entanto, são apenas algumas possibilidades, o \textbf{caminho} a ser trilhado necessita ser elaborado de forma a considerar a historicidade e os contextos dos sujeitos envolvidos, suas especificidades humanas.
Cabe aos educadores(as), em um processo contínuo de estudo e de formação continuada, buscarem o domínio desses conceitos³¹ e, como \textbf{autores} da sua própria prática pedagógica, galgarem os seus percursos a partir da realidade em que se encontram imersos.
Desse modo, desejamos que as escolas possam cada vez mais ser o lugar da cultura mais elaborada ( MELLO ; FARIAS, 2010 ) e que todos(as) os(as) educadores(as), possam enfrentar esse desafio, de \textbf{maneira} a vivenciarem, cotidianamente, experiências pedagógicas que possam garantir a todos(as) os(as) estudantes catarinenses o direito de aprender a ler e a escrever, de alfabetizar -se, no âmbito da alfabetização e do \textbf{letramento}, de modo indissociável e progressivamente, pleno.

% matched lemmas: abordagem, ancorar, análise, aprendiz, articulado, aula, autor, autêntico, auxiliar, caminho, capacidade, complexo, compreensão, conceitual, configurar, constituinte, constitutivo, convir, crítico, decorrer, depender, devido, diagnóstico, discursivo, dizer, embasar, erro, esperar, evento, exigir, final, humanidade, implicar, indicação, inerente, instância, inter-relacionar, lembrar, letramento, maneira, materializar, metodológico, método, organizador, permear, reelaborar, reiterar, remeter, segar, sistematização, subsidiar, subsídio, tanto, tarefa, teórico, transversal, verificar, vista
\end{textsample}
